\noindent
Krátké shrnutí problému (general background knowledge): automatické „detektory AI“ (nástroje, které se snaží rozpoznat, zda text vytvořil člověk nebo model) mají v praxi významná omezení: vysoký podíl falešně pozitivních i falešně negativních výsledků, závislost na konkrétních datech/tréninku detektoru a snadná obcházení úpravami textu. Z tohoto důvodu by pedagogická praxe neměla stavět na detektorech jako na primárním nebo rozhodujícím nástroji při posuzování akademické integrity. (general background knowledge)

\bigskip

\noindent
Praktický přístup doporučený místo spoléhání se na detektory:
\begin{enumerate}
  \item Návrh úloh odolných vůči triviálnímu využití AI
    \begin{itemize}
      \item Preferujte procesně orientované úlohy: postupné milníky, pracovní verze, commit‑history, průběžné přihlášky a záznamy o postupu.
      \item Využívejte lokálně specifická nebo datově závislá zadání (lokální data, vlastní experimenty, pozorování terénu), která je těžké zcela vygenerovat „z gruntu“ externím AI.
      \item Zařazujte časově omezené in‑class writings nebo krátké testy na porozumění v průběhu výuky.
    \end{itemize}

  \item Požadavek na procesní důkazy při odevzdání
    \begin{itemize}
      \item Vyžadujte pracovní verze, poznámky, commity z repozitáře nebo screencasty editace jako součást odevzdání (viz šablony do sylabů).
      \item Do LMS přidejte pole deklarace: checkbox, krátký text s uvedením nástroje/verze/promptu a povinnou přílohu s pracovním materiálem (viz formulace v sekci \emph{Šablony do sylabů a vzory citací AI}).
    \end{itemize}

  \item Hodnocení procesu a transparentnost v rubrikách
    \begin{itemize}
      \item Zahrňte do rubrik explicitní kritéria vztahující se k použití AI a k procesnímu důkazu (např. Transparentnost použití AI 0–5; Originalita a vlastní přínos 0–10; Doklad procesu 0–5). Tato jednoduchá škála lze převzít a upravit podle typu úkolu (viz šablony do sylabů).
    \end{itemize}

  \item Alternativy k automatickému odsouzení: ověření porozumění
    \begin{itemize}
      \item Realizujte krátká ústní ověření nebo follow‑up pohovory k ověření postupu a porozumění (strukturou otázek zaměřených na proces).
      \item Požadujte reflexi: stručný text, kde student popíše rozhodnutí, zdroje a konkrétní kroky, které vedly k výsledku.
    \end{itemize}

  \item Formativní přístupy a prevence
    \begin{itemize}
      \item Využívejte formativní hodnocení a průběžnou zpětnou vazbu, aby studenti rozvíjeli dovednosti psaní a argumentace nezávisle na jednorázovém „výsledku“.
      \item Nabídněte workshopy o etickém použití AI, o tom jak deklarovat výstupy a jak dokumentovat proces (proaktivní transparentnost snižuje pokušení podvádět).
    \end{itemize}
\end{enumerate}

\bigskip

\noindent
Krátký praktický checklist pro vyučující (před řešením podezření)
\begin{enumerate}
  \item Zkontrolujte, zda student odevzdal požadované procesní důkazy (drafty, commity, screencast). (viz šablona pole deklarace)
  \item Projděte timeline práce v repozitáři nebo logy editací; vyžádejte krátké ústní doplnění či obhajobu postupu.
  \item Pokud je podezření, postupujte podle institucionálních pravidel o akademické integritě; nevyvozujte závěry pouze z automatických detektorů. (general background knowledge)
  \item Pokud chcete monitorovat rizika systematicky, navrhněte redesign zadání tak, aby bylo ověřování jednodušší a méně nákladné (milníky, reflexe, lokální data).
\end{enumerate}

\bigskip

\noindent
Závěr: místo spolehnutí se na „černé skříňky“ detektorů doporučujeme kombinaci didaktického designu (odolná zadání), povinné dokumentace procesu, transparentních rubrik a krátkých ověření porozumění. Tím se snižuje motivace ke zneužití AI a současně se posiluje učení a akademická integrita. (general background knowledge)